\chapter{Manifolds with rough boundary}
\label{chap:roughboundary}

Manifolds with rough boundary, after \cite[\S3.5]{GN26}: the charts take values in locally
convex subsets with dense interior of the model space, with the calculus of such sets developed
in \cite{LCA}. This includes manifolds with boundary and with corners. We prove that the boundary
is well defined, and treat open subsets, graphs, subspaces and the tangent bundle.

\section{Manifolds with rough boundary}
\label{sec:rough-boundary}

Following \cite[\S3.5]{GN26}, manifolds may be modeled on subsets of $E$ that are not open: a
subset $V \subseteq E$ is \emph{locally convex} if every point of $V$ has a convex neighborhood
in $V$, and it has \emph{dense interior} if $V \subseteq \overline{V^0}$. The $\Ck{k}$ calculus
of \cite{LCA} (\cite[\S1.4]{GN26}) works on such sets: $f\colon V \to F$ is $\Ck{1}$ if it is
continuous, $\Ck{1}$ on $V^0$, and $\dd(f|_{V^0})$ extends continuously to $V \times E$. Below,
$r \in \N_0 \cup \{\infty\}$ and ``$\Ck{r}$ on $V$'' means $\Ck{k}$ on $V$ for every $k \le r$.
Mathlib's charted spaces need open targets, so rough atlases are a separate structure, built on
Mathlib's partial homeomorphisms.

\begin{definition}[Rough charts and rough atlases]
  \label{def:rough-atlas}
  \lean{LocallyConvex.Manifolds.ContDiffLCUpTo, LocallyConvex.Manifolds.IsRoughChart,
    LocallyConvex.Manifolds.RoughAtlas}
  \leanok
  A \emph{rough $E$-chart} of a topological space $M$ is a homeomorphism
  $\varphi\colon U_\varphi \to V_\varphi$ from an open subset of $M$ onto a locally convex subset
  of $E$ with dense interior. A \emph{$\Ck{r}$ atlas of rough charts} is a set $\mathcal{A}$ of
  rough charts whose domains cover $M$ and whose transition maps
  $\psi \circ \varphi^{-1}\colon \varphi(U_\varphi \cap U_\psi) \to E$ are $\Ck{r}$; for $r = 0$
  one requires moreover that $\varphi(x) \in V_\varphi^0$ if and only if $\psi(x) \in V_\psi^0$
  \cite[Definition 3.5.1(a), p.~278]{GN26}.
\end{definition}

The domains $\varphi(U_\varphi \cap U_\psi)$ are relatively open in $V_\varphi$, hence again
locally convex with dense interior, so the $\Ck{r}$ condition makes sense
\cite[Remark 3.5.2, p.~279]{GN26}.

\begin{theorem}[The boundary is well defined]
  \label{thm:rough-boundary}
  \lean{LocallyConvex.Manifolds.ContDiffLC.mem_interior_of_rightInvOn,
    LocallyConvex.Manifolds.IsRoughChart.mem_interior_of_trans,
    LocallyConvex.Manifolds.RoughAtlas.mem_interior_iff, LocallyConvex.Manifolds.RoughAtlas.boundary,
    LocallyConvex.Manifolds.RoughAtlas.mem_boundary_iff}
  \leanok
  \uses{def:rough-atlas}
  Let $\varphi$, $\psi$ be charts of a $\Ck{r}$ atlas of rough charts at $x \in M$. Then
  $\varphi(x) \in V_\varphi^0$ if and only if $\psi(x) \in V_\psi^0$
  \cite[Lemma 3.5.6, pp.~279--280]{GN26}. Hence the \emph{boundary} $\partial M$, the set of
  points $x$ with $\varphi(x) \notin V_\varphi^0$ for some chart $\varphi$ at $x$
  \cite[Definition 3.5.1(c), p.~278]{GN26}, can be tested in any chart.
\end{theorem}

\begin{proof}
  \leanok
  For $r = 0$ this is an axiom. Let $r \ge 1$, $\kappa = \psi \circ \varphi^{-1}$,
  $y = \varphi(x) \in V_\varphi^0$, and suppose $z = \kappa(y)$ is not interior. Take an open
  $W \ni z$ such that $C = V_\psi \cap W$ is convex with nonempty interior. By Hahn--Banach there
  is $\lambda \in E'$ with $\lambda < \lambda(z)$ on $C^0$, hence $\lambda \le \lambda(z)$ on
  $C$. Pick $a \in C^0$ and $v = z - a$, so $\lambda(v) > 0$. The chain rule applied to
  $\kappa \circ \kappa^{-1} = \id$ at $z$ gives $u$ with $\dd\kappa(y, u) = v$. For $t > 0$ small,
  $\kappa(y + tu) \in C$, so $t^{-1}(\lambda(\kappa(y + tu)) - \lambda(z)) \le 0$; these
  quotients converge to $\lambda(v) > 0$, a contradiction.
\end{proof}

\begin{remark}[Infinite-dimensional subtleties]
  \label{rem:rough-subtleties}
  For $r = 0$ the axiom on interior points cannot be dropped in infinite dimensions: a closed
  ball of an infinite-dimensional Banach space is homeomorphic to the whole space
  \cite[Remark 3.5.9, p.~280]{GN26}. The proof above uses Hahn--Banach, hence the local
  convexity of $E$ and of $V_\psi$, but no completeness. The boundary $\partial M$ is only a
  subset of $M$; it need not be a submanifold (corners of $[0,1]^2$), and no structure on it is
  constructed.
\end{remark}

\begin{definition}[Manifolds with rough boundary]
  \label{def:rough-manifold}
  \lean{LocallyConvex.Manifolds.RoughAtlas.IsCompatible, LocallyConvex.Manifolds.RoughAtlas.maximal,
    LocallyConvex.Manifolds.RoughAtlas.IsMaximal, LocallyConvex.Manifolds.RoughAtlas.maximal_isMaximal,
    LocallyConvex.Manifolds.RoughAtlas.maximal_atlas_eq, LocallyConvex.Manifolds.RoughAtlas.boundary_maximal,
    LocallyConvex.Manifolds.RoughManifold}
  \leanok
  \uses{def:rough-atlas, thm:rough-boundary}
  A \emph{$\Ck{r}$ manifold with rough boundary} is a topological space with a maximal $\Ck{r}$
  atlas of rough charts \cite[Definition 3.5.1(b), p.~278]{GN26}. Every rough atlas is contained
  in exactly one maximal atlas, the set of rough charts compatible with it, and the boundary does
  not change when the atlas is replaced by its maximal atlas.
\end{definition}

\begin{proof}
  \leanok
  The transition map between two charts compatible with $\mathcal{A}$ factors, near each point,
  through a chart of $\mathcal{A}$, and is $\Ck{r}$ by the chain rule.
\end{proof}

\begin{definition}[$\Ck{s}$ maps]
  \label{def:rough-contmdiff}
  \lean{LocallyConvex.Manifolds.RoughAtlas.ContMDiff}
  \leanok
  \uses{def:rough-atlas}
  A map $f\colon M \to N$ between manifolds with rough boundary is $\Ck{s}$ if it is
  continuous and all its local representatives
  $\psi \circ f \circ \varphi^{-1}\colon \varphi(f^{-1}(U_\psi) \cap U_\varphi) \to F$ are $\Ck{s}$
  \cite[Definition 3.5.11, p.~281]{GN26}.
\end{definition}

\begin{proposition}[Composition and chart independence]
  \label{prop:rough-comp}
  \lean{LocallyConvex.Manifolds.IsRoughChart.contDiffLCUpTo_comp,
    LocallyConvex.Manifolds.RoughAtlas.ContMDiff.comp, LocallyConvex.Manifolds.RoughAtlas.contMDiff_id,
    LocallyConvex.Manifolds.RoughAtlas.contMDiff_maximal_iff}
  \leanok
  \uses{def:rough-contmdiff, def:rough-manifold}
  Compositions of $\Ck{s}$ maps are $\Ck{s}$ \cite[Proposition 3.5.12, p.~281]{GN26}. For
  $s \le r$, a map is $\Ck{s}$ for given atlases if and only if it is $\Ck{s}$ for their maximal
  atlases.
\end{proposition}

\begin{proof}
  \leanok
  Near a point, $\psi \circ (g \circ f) \circ \varphi^{-1}
  = (\psi \circ g \circ \kappa^{-1}) \circ (\kappa \circ f \circ \varphi^{-1})$ for a chart
  $\kappa$ at the image point, and the chain rule on locally convex subsets with dense interior
  applies, since the domain of $\psi \circ g \circ \kappa^{-1}$ is relatively open in
  $V_\kappa$. Chart independence follows by composing with the identity maps
  $(M, \mathcal{A}_{\max}) \to (M, \mathcal{A})$ and $(N, \mathcal{B}) \to (N, \mathcal{B}_{\max})$,
  which are $\Ck{s}$ by definition of the maximal atlases.
\end{proof}

\begin{proposition}[Examples]
  \label{prop:rough-examples}
  \lean{LocallyConvex.Manifolds.RoughAtlas.ofChartedSpace, LocallyConvex.Manifolds.RoughAtlas.boundary_ofChartedSpace,
    LocallyConvex.Manifolds.RoughAtlas.prod, LocallyConvex.Manifolds.RoughAtlas.boundary_prod,
    LocallyConvex.Manifolds.RoughAtlas.ofSubset, LocallyConvex.Manifolds.RoughAtlas.boundary_ofSubset,
    LocallyConvex.Manifolds.RoughAtlas.ofConvex, LocallyConvex.Manifolds.RoughAtlas.halfSpace,
    LocallyConvex.Manifolds.RoughAtlas.unitInterval, LocallyConvex.Manifolds.RoughAtlas.boundary_unitInterval}
  \leanok
  \uses{def:rough-atlas, thm:rough-boundary, def:mb-manifold}
  \begin{enumerate}
    \item A manifold modeled on $E$ is a $\Ck{r}$ manifold with rough boundary, for every $r$,
      with empty boundary \cite[Example 3.5.3(a), p.~279; Remark 3.5.7, p.~280]{GN26}.
    \item Products of $\Ck{r}$ manifolds with rough boundary, with the product charts, are
      $\Ck{r}$ manifolds with rough boundary \cite[Example 3.5.5, p.~279]{GN26}, and
      $\partial(M \times N) = \partial M \times N \cup M \times \partial N$.
    \item A nonempty locally convex subset $U \subseteq E$ with dense interior, with the single
      chart $U \hookrightarrow E$, is a $\Ck{r}$ manifold with rough boundary, whose boundary is
      $U \setminus U^0$ \cite[Example 3.5.3(b), p.~279; Remark 3.5.7, p.~280]{GN26}. This
      applies to convex sets with nonempty interior, such as closed half-spaces
      $\{\lambda \ge 0\}$, $\lambda \ne 0$, and $[0, 1]$, whose boundary is $\{0, 1\}$.
  \end{enumerate}
\end{proposition}

\begin{remark}
  \label{rem:rough-dense}
  \cite[Example 3.5.3(b)]{GN26} is stated for locally convex subsets with nonempty interior;
  the inclusion is a rough chart only if the interior is dense, which is what we assume. For
  convex sets the two conditions agree.
\end{remark}

\begin{proposition}[Open subsets and full submanifolds]
  \label{prop:rough-full-submanifold}
  \lean{LocallyConvex.Manifolds.RoughAtlas.restrict, LocallyConvex.Manifolds.RoughAtlas.boundary_restrict,
    LocallyConvex.Manifolds.RoughAtlas.IsFullSubmanifold, LocallyConvex.Manifolds.RoughAtlas.fullSubmanifold,
    LocallyConvex.Manifolds.RoughAtlas.boundary_fullSubmanifold,
    LocallyConvex.Manifolds.RoughAtlas.mem_interior_inter_preimage_iff}
  \leanok
  \uses{prop:rough-examples}
  A nonempty open subset $U$ of a $\Ck{r}$ manifold with rough boundary $M$ is one, with the
  restricted charts, and $\partial U = U \cap \partial M$ \cite[Example 3.5.4, p.~279]{GN26}. A
  nonempty subset $N \subseteq M$ such that every point of $N$ has a chart $\varphi$ with
  $\varphi(U_\varphi \cap N)$ locally convex with dense interior (a \emph{full submanifold}) is a
  $\Ck{r}$ manifold with rough boundary, with
  $\partial N = N \cap (\partial M \cup (N \setminus N^0))$
  \cite[Definition 3.5.19 and (3.24), p.~284]{GN26}.
\end{proposition}

\begin{proof}
  \leanok
  The targets of the restricted charts are relatively open in, or locally convex subsets with
  dense interior of, the targets of the charts of $M$, and the transition maps are restrictions of
  those of $M$. A point $\varphi(x)$ is interior to $\varphi(U_\varphi \cap N)$ if and only if it is
  interior to $V_\varphi$ and $x$ is interior to $N$, which gives the boundary and condition (A3)
  for $r = 0$.
\end{proof}

\begin{definition}[Submanifolds of manifolds with rough boundary]
  \label{def:rough-submanifold}
  \lean{LocallyConvex.Manifolds.IsAdaptedTo, LocallyConvex.Manifolds.subChart,
    LocallyConvex.Manifolds.RoughAtlas.IsSubmanifold, LocallyConvex.Manifolds.RoughAtlas.submanifold,
    LocallyConvex.Manifolds.IsLocallyConvexSubset.preimage_clm}
  \leanok
  \uses{prop:rough-full-submanifold}
  Let $F \subseteq E$ be a closed vector subspace. A subset $N \subseteq M$ is a submanifold
  modeled on $F$ if every point of $N$ lies in a chart $\varphi$ with
  $\varphi(U_\varphi \cap N) = V_\varphi \cap F$ and $V_\varphi \cap F$ with dense interior relative to
  $F$ (S1), and, for $r = 0$, adapted charts agree on boundary points (S2). The charts
  $\varphi|_{U_\varphi \cap N}$ then form a $\Ck{r}$ atlas of rough $F$-charts of $N$
  \cite[Definition 3.5.14 and Remark 3.5.16, pp.~282--283]{GN26}; the transition maps are
  corestricted to $F$ by \cite[Lemma 1.4.16]{GN26}.
\end{definition}

\begin{proposition}[Graphs]
  \label{prop:rough-graph}
  \lean{LocallyConvex.Manifolds.roughGraphChart, LocallyConvex.Manifolds.isAdaptedTo_roughGraphChart,
    LocallyConvex.Manifolds.RoughAtlas.isCompatible_roughGraphChart,
    LocallyConvex.Manifolds.RoughAtlas.exists_roughGraphChart,
    LocallyConvex.Manifolds.RoughAtlas.isSubmanifold_graph,
    LocallyConvex.Manifolds.RoughAtlas.graphSubmanifold,
    LocallyConvex.Manifolds.RoughAtlas.boundary_graphSubmanifold}
  \leanok
  \uses{def:rough-submanifold}
  Let $M$ be a $\Ck{r}$ manifold with rough boundary modeled on $E$, $N$ a $\Ck{r}$ manifold without
  boundary modeled on $F$, and $f \colon M \to N$ a $\Ck{r}$ map. Every point of the graph of $f$
  lies in a chart of $M \times N$ adapted to the graph and to $E \times \{0\}$; for $r \geq 1$ the
  graph is a submanifold of $M \times N$ modeled on $E \times \{0\}$, with boundary
  $\{(x, f(x)) : x \in \partial M\}$ \cite[Proposition 3.5.17, pp.~283--284]{GN26}.
\end{proposition}

\begin{proof}
  \leanok
  For charts $\varphi$ of $M$ and $\psi$ of $N$, the chart
  $\theta(x, y) = (\varphi(x), \psi(y) - \psi(f(x)))$ is $g \circ (\varphi \times \psi)$ for the
  $\Ck{r}$ diffeomorphism $g(z, w) = (z, w - h(z))$, $h = \psi \circ f \circ \varphi^{-1}$, so it is
  compatible with the product atlas. Its target is relatively open in
  $\varphi(U_\varphi \cap f^{-1}(U_\psi)) \times F$, because $\psi$ has open target, and it maps the
  graph onto $E \times \{0\}$. For $r = 0$, G--N obtain a submanifold with prescribed boundary
  \cite[3.5.15]{GN26}, which we do not formalize.
\end{proof}

\begin{definition}[Manifolds with corners]
  \label{def:corners}
  \lean{LocallyConvex.Manifolds.quadrant, LocallyConvex.Manifolds.RoughAtlas.IsWithCorners,
    LocallyConvex.Manifolds.RoughAtlas.IsWithBoundary, LocallyConvex.Manifolds.RoughAtlas.quadrantAtlas,
    LocallyConvex.Manifolds.RoughAtlas.boundary_quadrantAtlas,
    LocallyConvex.Manifolds.RoughAtlas.isWithCorners_quadrantAtlas}
  \leanok
  \uses{prop:rough-examples}
  A manifold with rough boundary modeled on $\R^n$ is a manifold with corners (with boundary) if
  its maximal atlas contains charts around every point whose ranges are open subsets of
  $[0, \infty[^n$ (of $\R^{n-1} \times [0, \infty[$) \cite[Remark 3.5.10, p.~281]{GN26}. The
  quadrant itself is a manifold with corners, whose boundary is the set of points with a vanishing
  coordinate.
\end{definition}

\begin{definition}[The tangent bundle of a manifold with rough boundary]
  \label{def:rough-tangent-bundle}
  \lean{LocallyConvex.Manifolds.RoughAtlas.tangentCoordChange,
    LocallyConvex.Manifolds.RoughAtlas.tangentCoordChange_comp,
    LocallyConvex.Manifolds.RoughAtlas.exists_clm_tangentCoordChange,
    LocallyConvex.Manifolds.RoughAtlas.tangentCoordChange_inv,
    LocallyConvex.Manifolds.RoughAtlas.tangentBundleCore,
    LocallyConvex.Manifolds.RoughAtlas.TangentBundle, LocallyConvex.Manifolds.RoughAtlas.tangentChart}
  \leanok
  \uses{def:rough-atlas}
  Let $\mathcal{A}$ be a $\Ck{r}$ atlas of rough charts of $M$, $r \ge 1$. The tangent bundle
  $TM$ is glued from the sets $V_\varphi \times E$, $\varphi \in \mathcal{A}$, along the
  maps $T(\psi \circ \varphi^{-1})\colon (y, w) \mapsto (\psi\varphi^{-1}(y),
  \dd(\psi \circ \varphi^{-1})(y, w))$; its tangent charts are
  $T\varphi\colon TU_\varphi \to V_\varphi \times E$
  \cite[3.5.22, p.~285; Definitions 3.3.5 and 3.3.7, p.~257]{GN26}. In Lean, $TM$ is the total
  space of a fiber bundle core over the chart domains, with fiber $E$ and coordinate changes
  $\dd(\psi \circ \varphi^{-1})(\varphi(x), \cdot)$.
\end{definition}

\begin{remark}[Fibers over boundary points]
  \label{rem:rough-tangent-linear}
  At a boundary point $x$, $\dd(\psi \circ \varphi^{-1})(\varphi(x), \cdot)$ is the
  continuous extension of the differential from the interior, not a derivative at an interior
  point. It is nevertheless continuous linear \cite[Lemma 1.4.7]{GN26} and invertible, by the
  chain rule \cite[Proposition 1.4.10]{GN26} on the domains $\varphi(U_\varphi \cap U_\psi)$,
  which are locally convex with dense interior. Hence each fiber, also over $\partial M$, is a
  topological vector space isomorphic to $E$ through every chart; it is all of $E$, not a cone of
  inward vectors.
\end{remark}

\begin{theorem}[$TM$ is a manifold with rough boundary]
  \label{thm:rough-tangent-atlas}
  \lean{LocallyConvex.Manifolds.ContDiffLCUpTo.tangent,
    LocallyConvex.Manifolds.RoughAtlas.isRoughChart_tangentChart,
    LocallyConvex.Manifolds.RoughAtlas.tangentChart_trans_eqOn,
    LocallyConvex.Manifolds.RoughAtlas.tangentAtlas,
    LocallyConvex.Manifolds.RoughAtlas.boundary_tangentAtlas,
    LocallyConvex.Manifolds.RoughAtlas.contMDiff_proj}
  \leanok
  \uses{def:rough-tangent-bundle, thm:rough-boundary, def:rough-contmdiff}
  For $r \ge 1$, the tangent charts form a $\Ck{r-1}$ atlas of rough $E \times E$-charts of
  $TM$, with transition maps $T(\psi \circ \varphi^{-1})$; the boundary of $TM$ is
  $\pi^{-1}(\partial M)$, and $\pi\colon TM \to M$ is $\Ck{r-1}$
  \cite[3.5.22, p.~285; Lemma 3.3.8, p.~258]{GN26}.
\end{theorem}

\begin{proof}
  \leanok
  The targets $V_\varphi \times E$ are locally convex with dense interior. If $g$ is $\Ck{s}$ on
  $D$, $s \ge 1$, then $(y, w) \mapsto (g(y), \dd g(y, w))$ is $\Ck{s-1}$ on $D \times E$
  \cite[Definition 3.3.3, p.~256]{GN26}, since $\dd g$ is $\Ck{s-1}$ and pairs of $\Ck{s-1}$
  maps are $\Ck{s-1}$. As $(V_\varphi \times E)^0 = V_\varphi^0 \times E$, interior points of
  $TM$ are the vectors over interior points of $M$, which gives the boundary of $TM$ and, for
  $r = 1$, the axiom on interior points.
\end{proof}

\begin{theorem}[Tangent maps and the chain rule]
  \label{thm:rough-tangent-map}
  \lean{LocallyConvex.Manifolds.RoughAtlas.tangentMap,
    LocallyConvex.Manifolds.RoughAtlas.dLC_localRep_change,
    LocallyConvex.Manifolds.RoughAtlas.tangentChart_tangentMap,
    LocallyConvex.Manifolds.RoughAtlas.continuous_tangentMap,
    LocallyConvex.Manifolds.RoughAtlas.contMDiff_tangentMap,
    LocallyConvex.Manifolds.RoughAtlas.tangentMap_id,
    LocallyConvex.Manifolds.RoughAtlas.tangentMap_comp,
    LocallyConvex.Manifolds.RoughAtlas.tangentCoordChange_ofChartedSpace}
  \leanok
  \uses{thm:rough-tangent-atlas, prop:rough-comp}
  Let $f\colon M \to N$ be $\Ck{s}$, $s \ge 1$. The tangent map $Tf$, given in tangent charts by
  $T(\psi \circ f \circ \varphi^{-1})$, does not depend on the charts and is $\Ck{s-1}$; moreover
  $T(\id) = \id$ and $T(g \circ f) = Tg \circ Tf$ for $\Ck{1}$ maps
  \cite[Definition 3.3.11, p.~259; Proposition 3.3.12, p.~259; Remark 3.3.13, p.~260;
  3.5.22, p.~285]{GN26}. For ordinary manifolds the coordinate changes of $TM$ are those of
  Definition~\ref{def:tangent-bundle}.
\end{theorem}

\begin{proof}
  \leanok
  For charts $\varphi, \varphi'$ at $x$ and $\psi, \psi'$ at $f(x)$, both
  $\dd(\psi' \circ f \circ \varphi'^{-1})(\varphi'(x), \dd(\varphi' \circ \varphi^{-1})(\varphi(x), v))$
  and $\dd(\psi' \circ \psi^{-1})(\psi(f(x)), \dd(\psi \circ f \circ \varphi^{-1})(\varphi(x), v))$
  are equal to $\dd(\psi' \circ f \circ \varphi^{-1})(\varphi(x), v)$, by the chain rule on the
  common domain $\varphi(U_\varphi \cap U_{\varphi'} \cap f^{-1}(U_\psi \cap U_{\psi'}))$, which
  has dense interior; the differential on a smaller domain with dense interior is the
  restriction. The chain rule for $T(g \circ f)$ is the same computation. We do not claim this
  argument is new.
\end{proof}
